\documentclass[aspectratio=169,11pt]{beamer}

% ============================================
% PACKAGES
% ============================================
\usepackage[utf8]{inputenc}
\usepackage[T1]{fontenc}
\usepackage[ngerman]{babel}
\usepackage{lmodern}
\usepackage{tcolorbox}
\tcbuselibrary{skins, hooks}
\usepackage{listings}
\usepackage{xcolor}
\usepackage{fontawesome5}
\usepackage{tikz}
\usetikzlibrary{trees, shapes, arrows.meta}

% Modern Font
\IfFileExists{FiraSans.sty}{\usepackage[sfdefault]{FiraSans}}{}
\renewcommand{\familydefault}{\sfdefault}

% ============================================
% COLORS & DESIGN SYSTEM
% ============================================
\definecolor{BgColor}{HTML}{FAFAFA}
\definecolor{Primary}{HTML}{2B3A42}
\definecolor{Accent}{HTML}{E84A5F}
\definecolor{Secondary}{HTML}{3F5765}
\definecolor{CodeBg}{HTML}{F0F0F0}
\definecolor{Success}{HTML}{28A745}

\setbeamercolor{background canvas}{bg=BgColor}
\setbeamercolor{title}{fg=Primary}
\setbeamercolor{frametitle}{bg=Primary, fg=white}
\setbeamercolor{structure}{fg=Accent}
\setbeamercolor{normal text}{fg=Primary}
\setbeamercolor{itemize item}{fg=Accent}

\setbeamertemplate{navigation symbols}{}
\setbeamertemplate{footline}{%
  \leavevmode%
  \hbox{%
  \begin{beamercolorbox}[wd=.333333\paperwidth,ht=2.25ex,dp=1ex,center]{author in head/foot}%
    \usebeamerfont{author in head/foot}\tiny Falko Himmel
  \end{beamercolorbox}%
  \begin{beamercolorbox}[wd=.333333\paperwidth,ht=2.25ex,dp=1ex,center]{title in head/foot}%
    \usebeamerfont{title in head/foot}\tiny Grundlagen der Praktischen Informatik
  \end{beamercolorbox}%
  \begin{beamercolorbox}[wd=.333333\paperwidth,ht=2.25ex,dp=1ex,right]{date in head/foot}%
    \usebeamerfont{date in head/foot}\insertframenumber{} / \inserttotalframenumber\hspace*{2ex} 
  \end{beamercolorbox}}%
  \vskip0pt%
}

% ============================================
% CUSTOM BOXES
% ============================================
\tcbset{
    modernbox/.style={
        enhanced,
        boxrule=0pt,
        frame hidden,
        colback=white,
        coltitle=white,
        colbacktitle=Secondary,
        fonttitle=\bfseries,
        drop lifted shadow,
        arc=4pt,
        toptitle=2mm,
        bottomtitle=2mm,
        left=2mm,
        right=2mm
    },
    alertbox/.style={
        modernbox,
        colbacktitle=Accent
    }
}

% ============================================
% CODE LISTINGS
% ============================================
\lstdefinelanguage{Haskell}{
  keywords={class, data, deriving, do, else, if, in, import, let, module, of, then, type, where, case, error},
  keywordstyle=\color{Accent}\bfseries,
  identifierstyle=\color{Primary},
  sensitive=true,
  comment=[l]{--},
  commentstyle=\color{gray}\ttfamily\itshape,
  stringstyle=\color{teal}\ttfamily,
  morestring=[b]"
}

\lstset{
    language=Haskell,
    backgroundcolor=\color{CodeBg},
    basicstyle=\ttfamily\scriptsize,
    breaklines=true,
    frame=single,
    rulecolor=\color{CodeBg},
    frameround=tttt,
    numbers=left,
    numberstyle=\tiny\color{gray},
    xleftmargin=15pt,
    tabsize=4,
    showstringspaces=false
}

% ============================================
% TITLE PAGE
% ============================================
\title{\textbf{Tutorium 12}}
\subtitle{Sortieralgorithmen}
\institute{Grundlagen der Praktischen Informatik}
\author{}
\date{\today}

\begin{document}

\begin{frame}[plain]
    \titlepage
\end{frame}

\begin{frame}{Inhaltsverzeichnis}
    \tableofcontents
\end{frame}

% ============================================
% O(n^2) ALGOS
% ============================================
\section{Quadratische Sortieralgorithmen \texorpdfstring{$\mathcal{O}(n^2)$}{O(n\textasciicircum 2)}}

\begin{frame}{Selection Sort \& Bubble Sort}
    Beide Algorithmen haben eine worst-case Laufzeit von $\mathcal{O}(n^2)$.
    
    \begin{columns}[T]
        \begin{column}{0.48\textwidth}
            \begin{tcolorbox}[modernbox, title=Selection Sort]
                \begin{enumerate}
                    \item Finde das kleinste Element.
                    \item Entferne es aus der Liste.
                    \item Füge es an die sortierte Teilliste an.
                    \item Wiederholen.
                \end{enumerate}
            \end{tcolorbox}
        \end{column}
        \begin{column}{0.48\textwidth}
            \begin{tcolorbox}[modernbox, title=Bubble Sort]
                \begin{enumerate}
                    \item Gehe durch die Liste.
                    \item Vergleiche benachbarte Elemente.
                    \item Tausche sie, falls sie falsch herum stehen.
                    \item Wiederholen bis sortiert.
                \end{enumerate}
            \end{tcolorbox}
        \end{column}
    \end{columns}
\end{frame}

\begin{frame}{Bubble Sort (Visuell)}
    \begin{center}
    \begin{tikzpicture}[node distance=1.5cm, every node/.style={draw, rectangle, minimum size=0.8cm, font=\large}]
        \node (A) {5};
        \node (B) [right of=A, fill=Accent!20] {9};
        \node (C) [right of=B, fill=Accent!20] {2};
        \node (D) [right of=C] {1};
        \node (E) [right of=D] {6};
        
        \draw[->, thick, Accent] (B.north) to[bend left] node[above, draw=none] {\tiny Tausch!} (C.north);
    \end{tikzpicture}
    
    \vspace{0.5cm}
    \large Bei jedem Durchlauf steigen die größten Elemente wie ,,Blasen`` nach oben.
    \end{center}
\end{frame}


\begin{frame}{Insertion Sort}
    Ebenfalls $\mathcal{O}(n^2)$. Sehr intuitiv (ähnlich wie man Spielkarten sortiert).
    
    \begin{tcolorbox}[modernbox, title=Ablauf]
        \begin{enumerate}
            \item Nimm das nächste Element aus der unsortierten Liste.
            \item Finde die \textbf{korrekte Stelle} in der bereits sortierten Folge $S$.
            \item Füge es in $S$ ein.
            \item Wiederholen.
        \end{enumerate}
    \end{tcolorbox}
    
    \textbf{Beispiel ($S$ | $F$):} \\
    $() \quad | \quad (21, 5, 3)$ \\
    $(21) \quad | \quad (5, 3)$ \\
    $(5, 21) \quad | \quad (3)$ \\
    $(3, 5, 21) \quad | \quad ()$
\end{frame}

% ============================================
% O(n log n) ALGOS
% ============================================
\section{Log-Lineare Sortieralgorithmen \texorpdfstring{$\mathcal{O}(n \log n)$}{O(n log n)}}

\begin{frame}{Merge Sort}
    Ein klassischer Divide-and-Conquer-Algorithmus mit garantierter Laufzeit $\mathcal{O}(n \log n)$.
    
    \begin{tcolorbox}[modernbox, title=Ablauf]
        \begin{enumerate}
            \item \textbf{Teilen:} Teile die Liste in der Mitte in zwei Hälften $F_1, F_2$.
            \item \textbf{Rekursion:} Führe Merge Sort auf $F_1$ und $F_2$ aus.
            \item \textbf{Mergen:} Füge die sortierten Hälften wieder (in linearer Zeit) zusammen.
        \end{enumerate}
    \end{tcolorbox}
    
    
    \begin{center}
    \begin{tikzpicture}[level distance=1.2cm, sibling distance=2.5cm, every node/.style={rectangle, draw, minimum width=1.2cm, minimum height=0.6cm}]
    \node {[21, 5, 3, 17]}
      child {node {[21, 5]}
        child {node {[21]}}
        child {node {[5]}}
      }
      child {node {[3, 17]}
        child {node {[3]}}
        child {node {[17]}}
      };
    \end{tikzpicture}
    \end{center}
$
\end{frame}

\begin{frame}{Quick Sort}
    Sehr schnell in der Praxis, aber worst-case $\mathcal{O}(n^2)$. Average-case $\mathcal{O}(n \log n)$.
    
    \begin{tcolorbox}[modernbox, title=Ablauf]
        \begin{enumerate}
            \item \textbf{Pivot wählen:} Wähle ein Pivotelement.
            \item \textbf{Partitionieren:} Teile in $F_1$ (kleiner als Pivot) und $F_2$ (größer/gleich).
            \item \textbf{Rekursion:} Sortiere $F_1$ und $F_2$.
            \item \textbf{Zusammensetzen:} $F_1 ++ [Pivot] ++ F_2$
        \end{enumerate}
    \end{tcolorbox}
\end{frame}

% ============================================
% LIVE DEMO
% ============================================
\section{Praxis}

\begin{frame}[plain]
    \begin{center}
        \Huge \textbf{Live Demo} \\
        \vspace{0.5cm}
        \large \faPen~ \texttt{Altklausuraufgaben zu Sortieralgorithmen} \\
        \vspace{0.3cm}
        \normalsize Wir spielen die Algorithmen live an Beispielen durch!
\vspace{0.5cm}\begin{center}\href{https://moodle.uni-ulm.de/pluginfile.php/1348606/mod_folder/content/0/2024_2_exam.pdf}{\beamerbutton{2024\_2}}\quad\href{https://moodle.uni-ulm.de/pluginfile.php/1348606/mod_folder/content/0/2024_1_exam.pdf}{\beamerbutton{2024\_1}}\quad\href{https://moodle.uni-ulm.de/pluginfile.php/1348606/mod_folder/content/0/2023_2_exam.pdf}{\beamerbutton{2023\_2}}\quad\href{https://moodle.uni-ulm.de/pluginfile.php/1348606/mod_folder/content/0/2023_1_exam.pdf}{\beamerbutton{2023\_1}}\end{center}
    \end{center}
\end{frame}

\end{document}
